\subsubsection{Perceptrón multicapa}

Hemos usado el código de extracción de características para procesar los audios que tenemos en el dataset personalizado con 100 audios en español (50 reales y 50 generados a partir de esos reales). Una vez preprocesados los audios con sus características extraídas, empleamos una CNN (red neuronal convolucional) para sacar un embedding de 128 dimensiones del espectrograma del audio y un MLP (perceptrón multicapa) para sacar un embedding de 64 dimensiones de las características paralingüísticas del audio. Tras esto, las concatenamos teniendo un embedding de 192 dimensiones siendo las primeras 64 las características y las otras 128 las del espectrograma.

Después, creamos un perceptrón multicapa cuyo objetivo es clasificar el audio en si es generado o es real empleando para ello lógica difusa (\emph{fuzzy logic}). Para ello, pasamos la salida del perceptrón por un controlador fuzzy que irá aprendiendo también sus parámetros y adaptará la salida para que esté entre 0 y 1 siendo que cuanto más cercano al 0, más posible es que el audio sea real y no sea generado.
En la práctica, como el dataset es de 100 datos, no hay suficientes para entrenar un perceptrón de forma correcta ya que cuando hicimos las pruebas, todas las predicciones del perceptrón no se alejan del 0.5 que en este contexto de probabilidad continuo significa que el modelo está indeciso y no lo sabe clasificar como real o como generado.